\chapter{Symmetry in Physics}
\label{symmetry}
Symmetry enters classical and quantum physics at several logically distinct
levels, and it is a common source of confusion to treat them as one. The
symmetry group of a Lagrangian (or action) is, in general, only a
\emph{subgroup} of the symmetry group of the equations of motion (EOM) that it
generates; and the symmetry group respected by a \emph{particular} solution is,
in general, only a subgroup of the symmetry group of the EOM themselves, once
initial or boundary conditions are imposed. This paper gives a systematic
account of this three-tier hierarchy $G_L \subseteq G_{EOM}$ and
$G_{sol}\subseteq G_{EOM}$, illustrates each inclusion with worked
examples (hidden/dynamical symmetries, integrable systems, electromagnetic
duality; central-force orbits, diffusion with asymmetric data), and then turns
to the qualitatively different phenomenon of \emph{spontaneous symmetry
breaking} (SSB), in which the narrowing from $G_{EOM}$ down to the symmetry of
the physically realized state occurs \emph{without} any explicit asymmetric
input in the Lagrangian, Hamiltonian, or boundary data. We develop the
Goldstone-theorem structure of continuous SSB, its discrete analogue, the
Landau free-energy viewpoint, the subtleties of gauge-symmetry ``breaking''
(Elitzur's theorem, unitary gauge, the Higgs mechanism), and the role of the
thermodynamic/infinite-volume limit in making SSB robust in a way that ordinary
initial-condition selection is not. A unifying table and closing synthesis
relate all three phenomena within a single logical skeleton.

\section{Introduction}

Symmetry is often introduced to students as a single, monolithic idea: ``the
Lagrangian is invariant under such-and-such transformation, therefore there is
a conserved quantity, therefore the system respects the symmetry.'' In
practice, the situation is considerably richer, and the richness is not a
technicality but the source of some of the deepest phenomena in physics, from
the accidental degeneracies of the Kepler problem to the existence of the
electron mass and the pion.

It is useful to keep three distinct objects in view:
\begin{enumerate}
\item the \textbf{Lagrangian} (or action, or Hamiltonian) that \emph{defines} the dynamics;
\item the \textbf{equations of motion} (EOM) that it \emph{generates}; and
\item an actual \textbf{solution} of those equations, singled out by initial or boundary data, or, in a many-body/field-theoretic setting, the state that is actually realized as the ground state or equilibrium configuration of the system.
\end{enumerate}
Each of these objects has its own symmetry group, and the three groups need not
coincide. Two inclusions organize the entire subject:
\begin{itemize}
    \item $G_L \subseteq G_{EOM}$: \emph{every symmetry of the Lagrangian is a
        symmetry of the equations of motion, but the equations of motion can
        admit strictly more symmetry than is manifest in the Lagrangian one
        wrote down.} This is the subject of Section~\ref{sec:LvsEOM}.
    \item $G_{sol} \subseteq G_{EOM}$: \emph{a particular solution, or a
        particular physically realized state, generally respects only a
        subgroup of the symmetry of the equations that produced it.} There are
        two logically distinct mechanisms for this narrowing, and
        distinguishing them is the central conceptual payoff of this paper:
\begin{itemize}
\item \emph{Explicit (extrinsic) narrowing}, discussed in
    Section~\ref{sec:explicit}, in which the asymmetry is put in by hand
        through the choice of initial or boundary data. This is contingent:
        choosing symmetric data restores a symmetric solution, and averaging
        over the group orbit of solutions restores the full symmetry at the
        ensemble level.
\item \emph{Spontaneous narrowing}, discussed at length in
    Section~\ref{sec:SSB}, in which the Lagrangian, the equations of motion,
        \emph{and} the boundary conditions are all fully symmetric, and yet the
        stable, physically realized solution is not. This is not contingent: it
        is a structural consequence of the existence of a degenerate manifold
        of energy-minimizing (or free-energy-minimizing) configurations,
        together with an instability of the symmetric configuration and, in the
        case of a genuine phase transition, the thermodynamic (infinite-volume)
        limit.
\end{itemize}
\end{itemize}

Section~\ref{sec:synthesis} assembles the two inclusions and the two mechanisms
of narrowing into a single table and discusses precisely why spontaneous
symmetry breaking is \emph{not} simply a special case of ``asymmetric boundary
conditions,'' even though both phenomena produce a solution less symmetric than
the equations that govern it.

\section{Symmetry of the Lagrangian versus Symmetry of the Equations of Motion}
\label{sec:LvsEOM}

\subsection{Setup and definitions}

Let $q^a(t)$, $a=1,\dots,n$, be generalized coordinates on a configuration
space $Q$, with Lagrangian $L(q,\dot q,t)$ and action
\[
S[q] = \int_{t_1}^{t_2} L(q(t),\dot q(t),t)\,dt .
\]
The Euler--Lagrange (EL) operator acting on $L$ is
\[
[L]_a = \frac{\partial L}{\partial q^a} - \frac{d}{dt}\frac{\partial L}{\partial \dot q^a}
\]
and the equations of motion are $[L]_a=0$.

\begin{definition}[Variational / Noether symmetry]
An infinitesimal transformation $\delta q^a = \epsilon X^a(q,\dot q,t)$ is a \emph{variational symmetry} of $L$ if there exists a function $F(q,\dot q,t)$ such that, \emph{identically in $q(t)$} (not only on solutions), $ \delta L=\epsilon \tfrac{dF}{dt}$.
Equivalently, the action changes only by a boundary term, $\delta S = \epsilon\big[F\big]_{t_1}^{t_2}$, for arbitrary paths, not just extremal ones.
\end{definition}

\begin{definition}[Symmetry of the equations of motion]
A transformation $g$ (finite or infinitesimal) is a \emph{symmetry of the EOM}, or a \emph{dynamical symmetry}, if it maps every solution of $[L]_a=0$ to another solution of the \emph{same} equations, i.e. it maps the solution set to itself. No reference to the action or to boundary terms is required by this definition.
\end{definition}

\subsection{Noether's theorem and the forward inclusion}

\begin{proposition}
\label{prop:noether-forward}
Every variational symmetry of $L$ is a symmetry of the equations of motion. In
    symbols, $G_L \subseteq G_{EOM}$.
\end{proposition}

\begin{proof}[Proof (standard, physicist's level)]
Two facts are needed.
\emph{(a) The EL operator annihilates total time derivatives.} identically in $q(t)$ — not merely on shell.

\emph{(b) A change of path is a change of path, regardless of whether it is a symmetry.} 
    Consider the one-parameter family of paths $q_\epsilon(t) = q(t) + \epsilon X(q(t),\dot q(t),t) + O( \epsilon^2)$, for \emph{arbitrary} $q(t)$, not necessarily a solution. By hypothesis,
\[
L(q_\epsilon,\dot q_\epsilon,t) = L(q,\dot q,t) + \epsilon\, \frac{dF}{dt} + O(\epsilon^2)
\]
identically in $q(t)$. Hence the action evaluated on $q_\epsilon$ differs from
    the action evaluated on $q$ only by the boundary term
    $\epsilon[F]_{t_1}^{t_2}$, \emph{for every path, not only extremal ones}.
    Consequently $q(t)$ is a stationary point of $S$ under variations that
    vanish at the endpoints if and only if $q_\epsilon(t)$ is — because the two
    actions differ by a term that does not depend on the interior of the path
    at all. Since ``being a solution of $[L]_a=0$'' is exactly the statement of
    being such a stationary point, $q(t)$ solves the EOM iff $q_\epsilon(t)$
    does, to first order in $\epsilon$. Thus the transformation maps solutions
    to solutions: it is a symmetry of the EOM.
\end{proof}

The familiar conserved charge follows from the same computation kept on shell:
\[
Q = \frac{\partial L}{\partial \dot q^a} X^a - F, \qquad \frac{dQ}{dt}\Big|_{\text{on shell}} = 0 .
\]

\subsection{Why the inclusion can be strict: hidden and dynamical symmetries}
\label{sec:hidden}

The converse of Proposition~\ref{prop:noether-forward} fails in general: the
equations of motion routinely admit more symmetry than is manifest in the
Lagrangian one starts from. It is useful to separate \emph{manifest point
symmetries} of $L$ (transformations of $q$ and $t$ alone, geometrically visible
in the way $L$ is written), from \emph{generalized (dynamical) variational
symmetries}, which are allowed to depend on velocities or momenta and need not
correspond to any simple geometric transformation of configuration space, from
the full group $G_{EOM}g$ of symmetries of the equations of motion:
\[
G_L^{\text{point}} \;\subseteq\; G_L^{\text{gen}} \;\subseteq\; G_{EOM}g .
\]
It is a subtle question, belonging to the \emph{inverse problem of the calculus
of variations}, exactly when the second inclusion is an equality; under
favorable regularity assumptions many symmetries of a regular,
finite-dimensional Lagrangian system's EOM can indeed be exhibited as
\emph{generalized} variational symmetries (possibly of a Lagrangian differing
from the original by a boundary term), but this is not automatic, and in field
theory and in degenerate (constrained) systems it can fail outright. What
matters for the point of this paper is the \emph{first} inclusion: it is very
often strict, i.e., there is real, physically important symmetry of the EOM
that is simply invisible if one only stares at the naive, manifest symmetry of
$L$.

\begin{example}[Free particle: hidden $SL(2,mathbf R)$]
The Lagrangian $L=\tfrac12 m\dot x^2$ manifestly has the Galilei symmetry
    (translations, boosts) plus time translation. But the equation of motion
    $\ddot x=0$ has straight lines as its general solution, and the space of
    straight lines is left invariant by the larger, one-dimensional-time
    conformal (``Schr\"odinger/M\"obius'') group acting projectively on $t$,
\[
t \;\longmapsto\; \frac{\alpha t+\beta}{\gamma t+\delta}, \qquad \alpha\delta-\beta\gamma=1,
\]
together with a compensating rescaling of $x$. This $SL(2,mathbf R)$ ``dynamical''
    symmetry (closely related to the conformal/Schr\"odinger symmetry of the
    free Schr\"odinger equation, and to the Ermakov--Lewis invariant of the
    time-dependent oscillator) is not visible as a point symmetry of
    $L=\tfrac12 m\dot x^2$ written in the ordinary way, yet it is a genuine
    symmetry of $\ddot x=0$.
\end{example}

\begin{example}[Kepler problem: hidden $SO(4)$]
For $L = \tfrac12 m\dot{\bm r}^2 + k/r$, the manifest symmetry is $SO(3)$ (rotations) plus time translation. The equations of motion, however, admit the additional conserved Laplace--Runge--Lenz vector
\[
\bm A = \bm p \times \bm L - mk\,\hat{\bm r},
\]
whose associated transformations close, together with angular momentum, into an
    $SO(4)$ algebra (for bound orbits). This hidden symmetry is responsible for
    the non-precession of Kepler ellipses and for the accidental $n^2$
    degeneracy of the hydrogen spectrum; it is a symmetry of the dynamics, not
    a geometrically manifest symmetry of the Lagrangian's potential term.
\end{example}

\begin{example}[Integrable systems: toroidal symmetry from action-angle variables]
More generally, for a Liouville-integrable Hamiltonian system with $n$ degrees
    of freedom, one can pass to action-angle variables $(I_i,\theta^i)$ in
    which $\mathcal{H} = \mathcal{H}(I)$ depends only on the actions. The flow generated by
    $\mathcal{H}$ is then, trivially, an ($n$-torus of) simultaneous phase shifts
    $\theta^i \to \theta^i+c^i$, an abelian $U(1)^n$ symmetry group of the
    equations of motion. This symmetry is exact and completely general for
    integrable systems, yet it is essentially never visible as a point symmetry
    of the Lagrangian written in the original coordinates; the free-particle
    $SL(2,mathbf R)$ and the Kepler $SO(4)$ above are, in this sense, special
    cases of the general toroidal symmetry structure of integrable systems,
    made visible only after a nontrivial (often nonlocal, momentum-dependent)
    change of variables.
\end{example}

\begin{example}[Electromagnetic duality]
Source-free Maxwell's equations,
\[
\nabla\!\cdot\!\bm E=0,\quad \nabla\!\cdot\!\bm B=0,\quad \nabla\times \bm E=-\dot{\bm B},\quad \nabla\times\bm B=\dot{\bm E},
\]
are invariant under the duality rotation $(\bm E,\bm B)\to(\bm E\cos\theta+\bm
    B\sin\theta,\ -\bm E\sin\theta+\bm B\cos\theta)$. This is a genuine
    one-parameter (in fact $U(1)$) symmetry of the equations of motion. It is
    not a manifest symmetry of the standard vector-potential Lagrangian
    $\mathcal{L}=-\tfrac14 F_{\mu\nu}F^{\mu\nu}$, because a general duality-rotated
    field strength cannot be written as $\dd A$ for a single, ordinary gauge
    potential $A_\mu$ without enlarging the field content (e.g., introducing a
    second, dual potential, at the cost of manifest locality or manifest
    Lorentz covariance). Here the extra symmetry of the field equations is, in
    a precise sense, obstructed from being manifest at the level of the
    conventional action.
\end{example}

\begin{remark}[Noether's second theorem: the opposite phenomenon]
For completeness we note the converse-flavoured phenomenon that occurs for
    \emph{local} (gauge) symmetries, governed by Noether's \emph{second}
    theorem: an infinite-dimensional, spacetime-dependent symmetry of the
    Lagrangian (e.g., $A_\mu\to A_\mu+\partial_\mu\lambda(x)$) does not produce
    a conservation law of the usual kind but instead an off-shell
    identity/constraint among the equations of motion themselves (e.g. the
    Bianchi identity $\partial_{[\mu}F_{\nu\rho]}=0$, or the contracted Bianchi
    identity $\nabla^\mu G_{\mu\nu}\equiv0$ in general relativity). This shows
    that the relationship between Lagrangian symmetry and EOM structure is
    genuinely two-sided: global symmetries of $L$ can be smaller than
    symmetries of the EOM (this section's main theme), while local symmetries
    of $L$ manifest themselves as identities/constraints on the EOM rather than
    as extra conserved charges.
\end{remark}

\section{From Equations of Motion to a Particular Solution: Explicit Narrowing by Data}
\label{sec:explicit}

\subsection{The stabilizer of a solution}

Suppose $G_{EOM}g$ acts on the solution space $\mathcal S$ of $[L]_a=0$
(Section~\ref{sec:LvsEOM} guarantees this action is well defined at least for
$g\in G_L$, and by definition for all of $G_{EOM}g$). A single solution
$q(t)\in\mathcal S$, picked out by initial data $(q(t_0),\dot q(t_0))$ or by
boundary data, need not be a fixed point of this action. Its symmetry is the
\emph{stabilizer subgroup}
\[
    G_{sol}=\mathrm{Stab}_{G_{EOM}g}(q) = \{ g\in G_{EOM}g : g\cdot q = q \} \subseteq G_{EOM} .
\]
Generically $G_{sol} \subsetneq G_{EOM}g$: the group moves the given
solution to a \emph{different} member of $\mathcal S$ (still a valid solution —
the equations do not care — but a different trajectory), rather than mapping it
to itself.

\subsection{Examples}

\begin{example}[Central force problem]
The Kepler/central-force EOM has full $SO(3)$ symmetry
    (Section~\ref{sec:hidden} showed it is in fact larger still). A single
    orbit, however, has a definite, fixed angular momentum vector $\bm L = \bm
    r\times\bm p\neq 0$, determined by the initial position and velocity, and
    lies entirely in the plane orthogonal to $\bm L$. The stabilizer of this
    orbit inside $SO(3)$ is only the $SO(2)$ subgroup of rotations about the
    $\bm L$ axis (together with, for a closed orbit, possibly a discrete
    reflection). The other two rotational generators move the orbital plane to
    a different plane, i.e. map this solution to a genuinely different
    solution. Nothing is wrong with the dynamics; the reduction from $SO(3)$ to
    $SO(2)$ is entirely a consequence of the choice of initial data.
\end{example}

\begin{example}[Free particle]
$\ddot x=0$ is invariant under all spatial translations and Galilean boosts. A
    single trajectory $x(t)=x_0+v_0t$ is left invariant only by time
    translations (and, trivially, the identity among boosts/translations); a
    spatial translation or a boost maps it to a different straight line.
    Symmetric \emph{initial data} — e.g. $x_0=0,\,v_0=0$, the particle
    permanently at rest at the origin — is the unique choice whose stabilizer
    is the full translation-and-boost group; any other data breaks it
    explicitly.
\end{example}

\begin{example}[Diffusion equation]
The heat equation $\partial_t u = D\,\partial_x^2 u$ on $mathbf R$ is invariant
    under spatial translations, parity, and the parabolic scaling $x\to\lambda
    x,\ t\to\lambda^2 t$. An initial condition concentrated at the origin and
    even, $u(x,0)=\delta(x)$, yields a solution $u(x,t)$ that remains even in
    $x$ and self-similar under the parabolic scaling for all $t>0$ — the full
    symmetry \emph{compatible with an initial time slice} is preserved. A
    generic asymmetric initial condition (say, a delta function centered at
    $x=1$) produces a solution with no residual translation or parity symmetry
    at any later time. In both cases the equation itself is untouched; only the
    realized solution's symmetry depends on the data.
\end{example}

\subsection{The contingent character of this narrowing}

Two features distinguish this mechanism, and both will be denied in Section~\ref{sec:SSB}:
\begin{itemize}
\item \textbf{Removability.} Choosing symmetric data (when the symmetric
    configuration is dynamically consistent and stable) restores a symmetric
        solution. There is no obstruction internal to the dynamics.
\item \textbf{Restorability under averaging.} The full orbit $G_{EOM}g \cdot q
    = \{g\cdot q: g\in G_{EOM}g\}$ of all data-related solutions is, as a
        \emph{set} (or as a statistical ensemble/moduli space), fully
        $G_{EOM}$-symmetric. The asymmetry lives only in the choice of \emph{one
        representative}, not in the dynamics.
\end{itemize}

\section{Spontaneous Symmetry Breaking}
\label{sec:SSB}

\subsection{What is qualitatively different}

Spontaneous symmetry breaking (SSB) is often first described using the same
language as Section~\ref{sec:explicit} — ``the equations are symmetric but the
solution is not'' — and this can obscure the fact that SSB is mechanistically a
different phenomenon. In SSB, the asymmetry does \emph{not} need to be supplied
through the boundary or initial data: even fully symmetric data typically
\emph{fails} to select the symmetric configuration, because that configuration,
although it does solve the equations of motion and does respect the full
symmetry, is \emph{dynamically or thermodynamically unstable} — it is a local
maximum or saddle point of the energy (or free energy) rather than a minimum.
The system is driven — by an arbitrarily small perturbation, by quantum or
thermal fluctuations, or simply by the requirement of stability — onto one
member of a continuous or discrete \emph{degenerate family of asymmetric
minima}, and, in a system with infinitely many degrees of freedom, is unable to
tunnel or fluctuate back out of it.

\subsection{The prototype: complex scalar field with a Mexican-hat potential}

Consider a complex scalar field $\phi(x)$ with Lagrangian density
\[
\mathcal{L} = \partial_\mu\phi^*\partial^\mu\phi - V(|\phi|), \qquad
V(|\phi|) = -\mu^2|\phi|^2 + \lambda|\phi|^4, \qquad \mu^2,\lambda>0 .
\]
$\mathcal{L}$ (and the resulting field equations) is manifestly invariant under the
global phase rotation $\phi\to e^{i\alpha}\phi$, a $U(1)$ symmetry. Static,
spatially uniform, minimum-energy configurations minimize $V$; setting 
$dV/d|\phi|=0$ gives
\[
|\phi| = v \equiv \sqrt{\mu^2/2\lambda}, \qquad \arg\phi \ \text{arbitrary}.
\]
The minimum is degenerate along a circle $|\phi|=v$ — the \emph{vacuum
manifold} $\mathcal M \cong U(1)$. The symmetric configuration $\phi=0$ is a
solution of the field equations and is $U(1)$-invariant, but it sits at a
\emph{local maximum} of $V$ (since $\mu^2>0$), so it is unstable: any
infinitesimal perturbation grows and drives the field to the circle of minima.
Whichever point of the circle the system settles on, say $\phi = v$ (real,
$\theta_0=0$), that particular vacuum is left invariant by \emph{no} continuous
subgroup of $U(1)$ (other than the identity) — the symmetry is completely,
spontaneously broken.

\begin{definition}[Spontaneous symmetry breaking]
Let a theory have exact symmetry group $G$ (of the Lagrangian and hence the
    EOM). The theory exhibits SSB down to $H\subseteq G$ if the stable,
    physically realized ground state(s)/vacuum manifold $\mathcal M$ has
    stabilizer $H = \mathrm{Stab}_G(\phi_{\text{vac}}) \subsetneq G$ for a generic
    point $\phi_{\text{vac}}\in\mathcal M$, with the degenerate vacua
    parametrized by the coset space $G/H \cong \mathcal M$.
\end{definition}

This should be contrasted with \emph{explicit} breaking, in which a term that
is not $G$-invariant is added directly to $\mathcal{L}$ itself (e.g.
$\epsilon(\phi+\phi^*)$ above): there $G$ is not even a symmetry of the
equations of motion to begin with, the minimum is generically unique, and there
is nothing ``spontaneous'' about the resulting asymmetry.

\subsection{Goldstone's theorem}

Write $\phi = (v+\rho)\,e^{i\theta/v}$ near the chosen vacuum, with $\rho,\theta$ real fields. Expanding $V$,
\[
V = \mu^2\rho^2 + O(\rho^3), 
\]
with \emph{no} potential-energy dependence on $\theta$ at all (moving along the
vacuum circle costs zero energy, by construction, since $\mathcal M$ is a
manifold of degenerate minima). Consequently $\rho$ is a massive mode
($m_\rho^2\propto\mu^2$), while $\theta$ is an exactly massless mode: a
\emph{Nambu--Goldstone boson}. This is the content of Goldstone's theorem:

\begin{theorem}[Goldstone, 1961; Nambu, 1960]
When a continuous global symmetry group $G$ is spontaneously broken to a
    subgroup $H$, there appears one massless (Nambu--)Goldstone mode for each
    broken generator, i.e. $\dim(G/H)$ massless modes, corresponding to the
    flat directions of the potential tangent to the degenerate vacuum manifold
    $\mathcal M \cong G/H$.
\end{theorem}

\subsection{Examples of continuous SSB}

\begin{itemize}
\item \textbf{Ferromagnetism.} The Heisenberg Hamiltonian $\mathcal{H} =
    -J\sum_{\langle ij\rangle}\bm S_i\!\cdot\!\bm S_j$ ($J>0$) is exactly
        $O(3)$-invariant (global spin rotations). Below the Curie temperature
        the equilibrium state carries a spontaneous magnetization $\bm M\neq0$
        pointing in some (spontaneously chosen) direction, breaking $O(3)\to
        O(2)$. The Goldstone modes are spin waves (magnons), gapless in the
        isotropic limit.

\item \textbf{Crystallization.} A liquid's Hamiltonian is invariant under the
    full continuous Euclidean group (translations and rotations). A crystal's
        density is periodically modulated, breaking continuous translations
        down to a discrete lattice subgroup and continuous rotations down to
        the crystal's point group. The Goldstone modes are the acoustic phonon
        branches.

\item \textbf{Superfluidity / Bose--Einstein condensation.} The $U(1)$ symmetry
    associated with particle-number phase, $\hat\psi \to
        e^{i\alpha}\hat\psi$, is spontaneously broken by a nonzero condensate
        order parameter $\langle\hat\psi\rangle\neq0$. The Goldstone mode is
        the phonon (Bogoliubov) excitation.

\item \textbf{Chiral symmetry breaking in QCD.} In the limit of massless up and
    down quarks, the QCD Lagrangian has an (approximate) chiral $SU(2)_L\times
        SU(2)_R$ symmetry. The quark condensate $\langle\bar q q\rangle\neq0$
        spontaneously breaks this to the diagonal (vector) $SU(2)_V$ (isospin).
        The pions are the corresponding (pseudo-)Goldstone bosons; because the
        quark masses are small but not exactly zero, chiral symmetry is also
        \emph{explicitly} broken a little, giving the pions their small but
        nonzero mass — a clean illustration of explicit and spontaneous
        breaking operating together.
\end{itemize}

\subsection{Discrete symmetry breaking}

Not every symmetry is continuous. The Ising Hamiltonian $\mathcal{H}=-J\sum_{\langle
ij\rangle}s_is_j$ is invariant under the global $\mathbb Z_2$ symmetry $s_i\to
-s_i$ for all $i$. Below the critical temperature, two degenerate,
symmetry-related equilibrium states exist (magnetization up or down); the
system spontaneously selects one. Since $G/H$ is now a two-point discrete set
rather than a continuous manifold, there is no Goldstone boson (Goldstone's
theorem applies only to continuous symmetries); instead, the relevant
excitations are domain walls (kinks), topological defects interpolating between
the two degenerate vacua.

\subsection{The Landau free-energy viewpoint}

The same structure appears in Landau's mean-field theory of phase transitions.
One posits a free energy $F(m,T)$ for an order parameter $m$, constrained by
the symmetry $G$ of the underlying Hamiltonian to be a $G$-invariant function
of $m$ (e.g. $F(m)=F(-m)$ for Ising-type order parameters, or $F$ a function of
$|\phi|$ alone for a $U(1)$ order parameter). A standard truncated expansion,
\[
F(m,T) = a(T)\,m^2 + b\,m^4 + \cdots, \qquad b>0,
\]
with $a(T)$ changing sign at the critical temperature $T_c$, reproduces exactly
the structure of Section~4.2: for $a(T)>0$ ($T>T_c$) the unique, symmetric
minimum $m=0$ is stable; for $a(T)<0$ ($T<T_c$) it becomes unstable and the
true minima sit at $m=\pm\sqrt{-a/2b}\neq0$ (or on a circle/sphere for a
continuous order parameter), even though $F$ itself is fully $G$-invariant at
every temperature.

\subsection{Gauge symmetry: a genuine subtlety}

When the broken symmetry is \emph{local} (gauged) rather than global, the naive Goldstone-theorem story requires care.

\begin{itemize}
\item \textbf{Elitzur's theorem.} A strictly local gauge symmetry cannot be
    spontaneously broken in the operational sense above: in a properly
        gauge-fixed path integral (compact gauge group), the expectation value
        of any gauge-\emph{variant} local operator vanishes exactly, order by
        order, so there is no gauge-invariant sense in which
        ``$\langle\phi\rangle\neq0$'' by itself signals broken gauge symmetry.
        What is physical is not the value of $\phi$ in some gauge, but
        gauge-invariant data: the mass spectrum, gauge-invariant correlators,
        and the pattern of unbroken global symmetries.
\item \textbf{The Higgs mechanism.} The colloquial statement ``the gauge
    symmetry is spontaneously broken'' is best understood operationally through
        the \emph{unitary gauge} construction. Take the abelian Higgs model,
\[
\mathcal{L} = -\tfrac14 F_{\mu\nu}F^{\mu\nu} + |D_\mu\phi|^2 - V(|\phi|), \qquad D_\mu = \partial_\mu - i e A_\mu ,
\]
with $V$ as above. Writing $\phi=(v+\rho)e^{i\theta/v}$ and using the gauge
        freedom to set $\theta=0$ (unitary gauge), the covariant derivative
        term produces
\[
|D_\mu\phi|^2 \supset e^2v^2 A_\mu A^\mu ,
\]
a mass term for the gauge field, $m_A = ev$. The would-be Goldstone mode
        $\theta$ has disappeared from the physical spectrum; its single degree
        of freedom reappears as the new longitudinal polarization of the
        now-massive vector field (2 transverse $\to$ 3 polarizations), so the
        total count of propagating degrees of freedom is unchanged. The
        gauge-invariant statement of what happened is: \emph{the gauge boson
        acquires a mass and a longitudinal polarization, and the physical
        spectrum contains one massive scalar ($\rho$) instead of a massless
        scalar plus a massless gauge boson.}
\item \textbf{Electroweak symmetry breaking} follows the same pattern with
    $SU(2)_L\times U(1)_Y \to U(1)_{\text{em}}$: three of the four gauge
        generators are ``broken'' in this operational sense (their Goldstone
        partners are eaten by $W^\pm,Z^0$, which become massive), while the
        combination corresponding to electric charge remains an exact, unbroken
        gauge symmetry, leaving the photon massless.
\end{itemize}

\subsection{The necessity of the thermodynamic/infinite-volume limit}

For a system with \emph{finitely many} degrees of freedom, true SSB does not
occur, even when the potential has the double-well/Mexican-hat shape. A single
quantum particle in a symmetric double-well potential does not spontaneously
localize in one well: quantum tunneling mixes the two classically degenerate
configurations, and the true ground state is the symmetric (or antisymmetric)
combination, restoring the symmetry exactly. SSB in the strict sense requires
the \emph{thermodynamic limit} (infinite volume, or infinitely many degrees of
freedom): only then does the tunneling amplitude between distinct vacua vanish
(formally, the distinct vacua generate unitarily inequivalent representations
of the canonical (anti)commutation relations — distinct superselection
sectors), so that a system prepared in one vacuum stays there. This is also the
content of the statement, familiar from statistical mechanics, that a genuine
phase transition (a non-analyticity of the free energy) can only occur in the
thermodynamic limit.

Equivalently, in the language of ergodic theory: below the transition, the
dynamics undergoes \emph{ergodicity breaking} — the accessible phase space (or
Hilbert space) splits into disjoint sectors, each respecting only $H\subsetneq
G$, and time (or ensemble) averages taken within a single sector do not
reproduce the fully $G$-symmetric answer, even though the underlying dynamics
is exactly $G$-symmetric.

\section{Synthesis: Two Inclusions, Two Mechanisms of Narrowing}
The essential logical point, worth stating explicitly, is this: \textbf{one can
always choose perfectly symmetric initial or boundary data, and the symmetric
field configuration ($\phi=0$ in the scalar example, $m=0$ in the Landau
example, $\bm M=0$ in the ferromagnet) is genuinely a solution of the equations
of motion, respecting the full symmetry.} What distinguishes SSB from the
phenomenon of Section~\ref{sec:explicit} is that this symmetric solution is
\emph{dynamically unstable} (a local maximum or saddle of the energy/free
energy rather than a minimum), so it is not the state that is physically
realized or thermodynamically stable; the system is driven onto one of a
continuous or discrete family of degenerate, asymmetric, stable configurations,
and — crucially, only in the presence of infinitely many degrees of freedom —
cannot fluctuate or tunnel back to a symmetric superposition.
Section~\ref{sec:explicit}'s narrowing is about \emph{which member of an
already-degenerate family of equally good, equally stable solutions you
happened to pick by fixing data}; Section~\ref{sec:SSB}'s narrowing is about
the fact that \emph{the symmetric configuration was never a viable stable state
to begin with}, so no choice of symmetric data could have preserved it.

\section{Worked Appendix: Full Mass Spectrum Near a Broken Vacuum}

To make the Goldstone counting fully explicit, return to $V(\phi) =
-\mu^2|\phi|^2+\lambda|\phi|^4$ and write $\phi = \phi_1+i\phi_2$ in real
components. The Hessian of $V$ at the symmetric point $\phi=0$ is
\[
\frac{\partial^2 V}{\partial\phi_i\partial\phi_j}\Big|_{0} = -2\mu^2\,\delta_{ij},
\]
negative definite: $\phi=0$ is a maximum in every direction, confirming its
instability directly (not merely along one special direction). At the true
vacuum $\phi_1=v,\ \phi_2=0$,
\[
\frac{\partial^2V}{\partial\phi_1^2}\Big|_{\text{vac}} = 8\lambda v^2 = 4\mu^2 \;(\text{massive, } \rho \text{ direction}), \qquad
\frac{\partial^2V}{\partial\phi_2^2}\Big|_{\text{vac}} = 0 \;(\text{massless, Goldstone direction}),
\]
with vanishing off-diagonal term. This is the same content as the
polar-coordinate computation of Section~5.3, presented in Cartesian form:
exactly one flat (massless) direction for the one broken generator of the
one-dimensional group $U(1)$, consistent with $\dim(G/H) = \dim U(1) -
\dim\{e\} = 1$.

For a general theory with symmetry group $G$ spontaneously broken to $H$, the
same computation — the Hessian of the potential restricted to the tangent space
of the vacuum manifold $\mathcal M \cong G/H$ vanishing identically, because
$\mathcal M$ is by definition a manifold of degenerate minima — gives the
general Goldstone count $\dim G - \dim H$, matching the theorem of Section~5.3.

\section{Concluding Remarks}

The statement that ``symmetry of the equations of motion may be wider than that
of the Lagrangian, and symmetry of a solution may be narrower than that of the
equations, due to fixed initial and/or boundary conditions'' packages two
independent facts, both traced to first principles above: the forward direction
of Noether's theorem does not have an automatic converse at the level of
manifest, geometric symmetry (Section~\ref{sec:LvsEOM}), and any group action
on a solution space generically has nontrivial-but-proper stabilizers once a
particular orbit is singled out by data (Section~\ref{sec:explicit}).
Spontaneous symmetry breaking sharpens this second point into something much
stronger and more physically consequential: a narrowing of symmetry that
survives even the most symmetric possible choice of data, rooted instead in the
instability of the symmetric extremum and, for genuine phase transitions, in
the thermodynamic limit (Section~\ref{sec:SSB}). Keeping these three levels —
Lagrangian, equations of motion, realized state — and the two distinct
mechanisms of narrowing conceptually separate is what allows one to correctly
diagnose, in any given physical system, \emph{why} a symmetry that is
manifestly present in the dynamics fails to be manifest in what is actually
observed.

